\begin{proof}
\begin{enumerate}[label=\arabic*.]
\item Let $v\in\mathbb Z$, put $q=v^2-v+1$, and let
$Ae_j=e_{j+1}$ on $\mathbb Q^5$, with indices modulo five. Define
\[
 X=(e_0,\ e_2+ve_4),\qquad
 Y=\begin{pmatrix}
 0&-q&-v^2&-1&v\\
 0&1-vq&-vq&v^2-v&q
 \end{pmatrix},\qquad N=XY.
 \tag{F.U1}\label{maps:eq:F.U1}
\]
Multiplying $Y$ by the two columns of $X$ gives $YX=0$, so
$N^2=X(YX)Y=0$. Consequently
\[
 C=(I+N)A,\qquad B=A^4(I-N)=C^{-1}.
\]
The cycle gives $A^5=I$. Exact multiplication and determinant
expansion give
\[
 N^2=0,\qquad A^5=C^4=I_5,\qquad
 \chi_C(x)=(x-1)(x+1)^2(x^2+1).
 \tag{F.U2}\label{maps:eq:F.U2}
\]
For the determinant calculation one may substitute
\[
 (xI-A)^{-1}=
 \frac{x^4I+x^3A+x^2A^2+xA^3+A^4}{x^5-1}
\]
in the rank-two determinant identity
\[
 \det(xI-C)=\det(xI-A)
       \det\bigl(I_2-YA(xI-A)^{-1}X\bigr).
\]
The identity $C^4=I$ is checked separately; it is not inferred
from the repeated factor $x+1$ of the characteristic polynomial.
The factor $x^2+1$ shows that $B$ and $C$ have exact order four.

\item Let $P$ have columns $\mathbf1/q,e_1,e_2,e_3,e_4$, and put
$A'=P^{-1}AP$, $B'=P^{-1}BP$. Since $q>0$ and $\det P=1/q$,
these are actual rational changes of coordinates. Define
\[
 s=v^2+v-1,\qquad r=-\frac{(v+1)(v^2-2)}2,
 \qquad
 \beta=(s,\ s-vr,\ v^2s+vr,\ r,\ -vs-r)
 \tag{F.U3}\label{maps:eq:F.U3}
\]
and
\[
 \widehat A=\begin{pmatrix}1&0\\0&A'\end{pmatrix},\qquad
 \widehat B=\begin{pmatrix}1&\beta\\0&B'\end{pmatrix}.
 \tag{F.U4}\label{maps:eq:F.U4}
\]

The actual lattice $P\mathbb Z^5$ has basis
$f_0=\mathbf1/q,f_1=e_1,\ldots,f_4=e_4$.
Its images under $A$ are $f_0,f_2,f_3,f_4$ and
$qf_0-f_1-f_2-f_3-f_4$, respectively; $A^{-1}=A^4$
therefore also preserves this lattice. The columns of $X$ are
$qf_0-f_1-f_2-f_3-f_4$ and $f_2+vf_4$, while
\[
 Y':=YP=\begin{pmatrix}
 -2&-q&-v^2&-1&v\\
 2(1-v)&1-vq&-vq&v^2-v&q
 \end{pmatrix}.
\]
Thus $X$ and $Y$ are integral maps between the specified lattices.
The already proved $YX=0$ shows that $I+N$ and its inverse $I-N$
preserve $P\mathbb Z^5$. Since $B=A^{-1}(I-N)$, both
$A',B'$ belong to $\operatorname{GL}_5(\mathbb Z)$.
For $r$, odd $v$ makes $v+1$ even, and even $v$ makes
$v^2-2$ even; hence $r$ and all entries of $\beta$ are integers.
The hatted matrices are therefore integral automorphisms.

The exact unit needed in the hatted lattice is
\[
 z=(-v^3+v-1,0,0,2,2(1-v^2))^t.
\]
Writing $a=-v^3+v-1$, the first coordinate of $Y'z$ is
$-2a-2+2v(1-v^2)=0$. Half its second coordinate is
$(1-v)a+(v^2-v)+q(1-v^2)=0$, since the first two terms
sum to $v^4-v^3+v-1$ and the last is its negative. Furthermore
\[
 \beta z=s(v^3-v-1)+2v^2r=1.
 \tag{F.U5}\label{maps:eq:F.U5}
\]
Indeed the two summands expand to
$v^5+v^4-2v^3-2v^2+1$ and
$-v^5-v^4+2v^3+2v^2$.
Since $AB=I-N$ and $NPz=X(Y'z)=0$, the actual six-dimensional
map satisfies
\[
 (\widehat A\widehat B-I)(0,z)^t=e_0.
\]
This establishes the unit in the displayed hatted coordinates.

The row identity
$\beta(I+B'+B'^2+B'^3)=0$ proves
$\widehat B^4=I$; also $\widehat A^5=I$. Their quotient actions
show that these orders are exact.

Use the specified lattice change
\[
 L=\frac14\begin{pmatrix}
 -2&3&-5&-1&5&-3\\
 0&2&-2&2&-2&2\\
 0&-4&0&0&4&-4\\
 0&0&0&-4&4&-4\\
 0&0&4&-4&4&-4\\
 0&0&0&-4&4&0
 \end{pmatrix},\qquad
 A_v=L^{-1}\widehat A L,\quad B_v=L^{-1}\widehat B L.
 \tag{F.U6}\label{maps:eq:F.U6}
\]
Here $\det L=-1/4$. Expanding these products gives entries in
$\mathbb Z[v]$ for both matrices. The complete expanded entries
and the exact identities
$LA_v=\widehat A L$, $LB_v=\widehat B L$ are supplied by the
polynomial certificate described in the last item. Thus the
defined matrices act on the stated $\mathbb Z^6$ for every integer
parameter. Their determinants are one and their inverses are
the integer powers $A_v^4,B_v^3$. Their characteristic polynomials are
\[
 (x-1)^2\Phi_5(x),\qquad (x-1)^2(x+1)^2(x^2+1),
\]
respectively.

\item Put $K=B_vA_v-I$. Multiplication gives $K^2=0$.
The following are actual third minors of $K$, in zero-based indices:
\[
 \begin{array}{c|c|c}
 \text{rows}&\text{columns}&\text{determinant}\\\hline
 (0,1,4)&(1,2,3)&v(v+1)\\
 (0,1,4)&(1,3,4)&v^2\\
 (0,3,4)&(1,2,3)&(v+1)^2.
 \end{array}
\]
Their integer combination $(v+1)^2-2v(v+1)+v^2$ equals one.
Hence $K$ has rank at least three, while $K^2=0$ in dimension six
gives rank at most three. The gcd of its third minors is one,
so its three nonzero Smith factors are all one. This proves
primitivity directly in the specified lattice, for every $v$.
For $M_1=A_v^t$ and $M_2=B_v^t$ the actual cusp map is
\[
 M_0=(M_1M_2)^{-1}=(I+K^t)^{-1}=I-K^t.
\]
Thus $M_0-I=-K^t$, which has the same integral Smith factors
and the required maximal square-zero type.

\item Both $A_v,B_v$ fix $e_0$. The actual quotient morphism to
the original rank-five pair is
\[
 R=P\,\operatorname{pr}_5 L:\mathbb Q^6\longrightarrow\mathbb Q^5,
 \qquad RA_v=AR,\quad RB_v=BR.
\]
It is surjective and has kernel $\mathbb Qe_0$, since
$Le_0=-e_0/2$. We prove that this quotient is irreducible over
$\mathbb Q$. The coprime factors $x-1$ and $\Phi_5(x)$ of
$\chi_A$ give the projector
$(I+A+A^2+A^3+A^4)/5$. Its image is $\mathbb Q\mathbf1$ and
its kernel is $S=\{\sum x_i=0\}$. Irreducibility of $\Phi_5$
shows that these are the only nonzero proper rational
$A$-invariant subspaces.

Neither is $C$-invariant. Indeed,
\[
 N\mathbf1=-2qX\binom{1}{v-1}
\]
has nonzero first entry and zero entries in positions one and three,
so $C\mathbf1=\mathbf1+N\mathbf1$ is not a multiple of
$\mathbf1$. Write $\epsilon=(1,1,1,1,1)$. If $S$ were
$C$-invariant, then $\epsilon N$ would be a multiple of
$\epsilon$; its first entry is zero, so that multiple would
vanish. But its entries in columns three and four are
$v^3-v-1$ and $v^3+v+1$. Their simultaneous vanishing would
give $v=-1$, at which the first is $-1$. This is impossible.
Since $B=C^{-1}$, the pair $(A,B)$ is irreducible, as claimed.

\item The rank-two update gives exact traces
\[
 \operatorname{tr}(AC)=\operatorname{tr}(YA^2X)=1-v^2q,
 \qquad
 \operatorname{tr}((AC)^{-1})
       =-\operatorname{tr}(YA^3X)=1+v^2q.
\]
Here $\operatorname{tr}(A^2)=\operatorname{tr}(A^3)=0$ and
$B=A^4(I-N)$ give the displayed equalities. An equivariant
isomorphism $F$ to the contragredient representation would
satisfy $FA=A^{-t}F$ and $FB=B^{-t}F$, and therefore identify
the character of the word $AC=AB^{-1}$ with the character
of its contragredient matrix. The two displayed traces would
be equal. Their difference $-2v^2q$ is nonzero when $v\ne0$,
proving the required nonselfduality.

For completeness, at $v=0$ the symmetric matrix
$Q_{ij}=1$ for $j-i\equiv1,4\pmod5$ and $Q_{ij}=0$ otherwise
has determinant two and satisfies $A^tQA=Q$, $B_0^tQB_0=Q$.
It gives the actual equivariant isomorphism in that case.

For the rank-six matrices, the same word has trace
\[
 h(v):=\operatorname{tr}(A_vB_v^{-1})=2-v^2q,
 \qquad h(v+1)-h(v)=-4v^3-3v^2-3v-1.
\]
Thus $h$ is strictly decreasing on the nonnegative integers,
in particular on the nonnegative multiples of twenty.
A simultaneous rational conjugacy of the ordered matrices
would preserve this trace. There are consequently infinitely
many such rational conjugacy classes already in that subfamily.

\item For $v\in20\mathbb Z$, put
$a=(1,3,-2,-1,2,0)^t$, $e=(1,0,0,0,0,0)^t$, and define
\[
 w_1=\frac15\sum_{j=0}^4(A_v^t)^j a,
 \qquad
 w_2=-\frac14\sum_{j=0}^3(B_v^t)^j e.
 \tag{F.U7}\label{maps:eq:F.U7}
\]
The polynomial matrices reduce to their parameter-zero
values modulo five and four, respectively. At zero,
$M_1a=a$ and $M_2e=e$. The two sums are therefore divisible
by their indicated denominators, so $w_1,w_2$ are integral.
Their defining averages give $M_iw_i=w_i$. The common
first-coordinate covector $\ell$ is fixed by both $M_i$ and
satisfies $\ell(w_1)=1$, $\ell(w_2)=-1$.

For $(m_1,m_2)=(5,4)$ let
$g_i(x)=M_ix+w_i/m_i$ on $\mathbb R^6/\mathbb Z^6$.
Then $g_i^j(x)=M_i^jx+jw_i/m_i$, so $g_i^{m_i}$ is the
identity. Every power with $0<j<m_i$ changes the actual
$\ell$-circle coordinate by $j/5$ or $-j/4$, respectively.
It has no fixed point. These are free actions of the exact
orders five and four, with the specified affine lifts.

\item Substitution of $v=0$ in the displayed construction gives
exactly the marked pair in Theorem~\ref{maps:src:F.2.5}, and
the averages give $w_1(0)=a$, $w_2(0)=-e$. This identifies
both the matrices and the affine numerators without a further
change of basis.

The independent certificate
\texttt{F2\_16\_family\_fresh.py} records the complete expanded
polynomial matrices and checks every matrix, quotient, trace
and affine identity above over $\mathbb Q[v]$, with coefficient
integrality in $\mathbb Z[v]$. It checks all four hundred
third minors as well as the displayed integer combination.
Its separate Smith computations at ten stated integer
parameters are additional finite checks; the universal
primitivity assertion follows from the displayed identity
equal to one. The conclusions concern the unrestricted
input tuples in the statement.
\end{enumerate}
\end{proof}
